\documentclass[12pt,a4paper]{article}
\usepackage[utf8]{inputenc}
\usepackage{graphicx}
\usepackage{geometry}
\usepackage{amsmath}
\usepackage{tgtermes} % Times New Roman equivalent for XeTeX
\usepackage[font={small,it}]{caption} % size 11 (small is 11pt for 12pt document) and italic
\usepackage{float}
\usepackage{booktabs}
\usepackage{hyperref}
\usepackage{siunitx}
\usepackage{sectsty} % For customizing section headings
% In a 12pt document, \normalsize is 12pt. \bfseries makes it bold.
\sectionfont{\fontsize{12}{14}\bfseries\selectfont} 

\geometry{
  a4paper,
  left=25mm,
  right=25mm,
  top=25mm,
  bottom=25mm,
}

\begin{document}

% --- Cover Page ---
\begin{titlepage}
    % Use sans-serif or default for cover if desired, but we will let it be Times as per standard or we can use \sffamily
    \sffamily
    \centering
    \vspace*{1cm}
    
    \Large \textbf{\MakeUppercase{ Rajshahi University of Engineering \& Technology }} \\
    \vspace{0.5cm}
    \large \textbf{\MakeUppercase{ Department of Electrical and Electronic Engineering }} \\
    
    \vspace{2.5cm}
    \Large \textbf{LABORATORY REPORT} \\
    \vspace{1cm}
    
    \begin{tabular}{ll}
        \textbf{Course No:} & EEE 3102 \\
        \textbf{Course Title:} & Digital Signal Processing Sessional \\
    \end{tabular}
    
    \vspace{2cm}
    
    \textbf{Experiment No:} 09 \\
    \vspace{0.5cm}
    \textbf{Name of the Experiment:} \\
    \Large \textbf{ Test Experiment about triac and diac together } \\
    
    \vspace{2.5cm}
    
    \begin{minipage}[t]{0.4\textwidth}
        \begin{flushleft} \large
            \emph{Submitted by:}\\
            John Doe \\
            Roll: 1901000 \\
            Section: A, Group: 1
        \end{flushleft}
    \end{minipage}
    \hfill
    \begin{minipage}[t]{0.5\textwidth}
        \begin{flushright} \large
            \emph{Submitted to:} \\
            
            Dr. Jane Smith \\
            Professor \\
            \vspace{0.3cm}
            
            Mr. Bob Brown \\
            Lecturer \\
            
            
        \end{flushright}
    \end{minipage}

    \vfill
    \textbf{Date of Performance:} October 09, 2026 \\
    \textbf{Date of Submission:} October 16, 2026
    
\end{titlepage}

% --- Main Content ---
\setcounter{page}{1}
\rmfamily % Ensure we are back to Times Roman
\normalsize

\begin{center}
    \textbf{Experiment No:} 09 \\
    \vspace{0.2cm}
    \textbf{Name of the Experiment:} Test Experiment about triac and diac together
\end{center}
\vspace{0.5cm}


\section{Objectives}
\begin{itemize}

    \item To demonstrate the bidirectional conduction characteristics of a TRIAC when triggered by a gate signal.

    \item To investigate the breakover voltage and bidirectional switching behavior of a DIAC.

    \item To observe and measure the AC output voltage across a load resistor controlled by a DIAC-TRiac combination circuit.

    \item To analyze the relationship between the gate trigger timing and the resulting phase-controlled AC output.

\end{itemize}


\section{Introduction}
The TRIAC, or Triode for Alternating Current, is a solid-state switch that controls power in either direction through the load. It essentially consists of two SCRs connected in antiparallel, allowing it to conduct current in both positive and negative half-cycles of an AC waveform. A TRIAC remains in a blocking state until a specific gate trigger voltage is applied; once triggered, it latches into the conduction state until the current drops below the holding current. This device is fundamental in AC power control applications, such as light dimmers and motor speed controllers, because it can be used to alter the phase angle of the output waveform.

The DIAC, or Diode for Alternating Current, is a two-layer, two-terminal semiconductor device that acts as a bidirectional trigger switch. Unlike a standard diode, a DIAC has no gate terminal and does not have a distinct anode or cathode. It remains in a high-impedance blocking state until the voltage across its terminals exceeds a specific breakover voltage, at which point it breaks down and conducts heavily in both directions. In a typical control circuit, the DIAC is placed in series with the gate of a TRIAC. The DIAC provides a consistent, sharp trigger pulse to the TRIAC's gate, ensuring reliable latching of the TRIAC at specific points in the AC cycle. Together, these components form the core of many AC phase-control systems.


\begin{figure}[H]
    \centering
    \includegraphics[width=0.8\textwidth]{/home/sword/Documents/LAB_Expert/labgen/runs/test_experiment_about_triac_and_diac_together/figs/theory_img.png}
    \caption{Reference diagram}
\end{figure}


\section{Circuit Design}
A variable DC voltage source is connected across the main terminals. A gate current is provided to trigger the device. The voltage is swept from negative to positive values.


\begin{figure}[H]
    \centering
    \includegraphics[width=0.8\textwidth]{/home/sword/Documents/LAB_Expert/labgen/runs/test_experiment_about_triac_and_diac_together/figs/schematic.png}
    \caption{Experimental Circuit Diagram}
\end{figure}


\section{Required Apparatus}
\begin{itemize}

    \item DC Power Supply (Variable) (Quantity: 1)

    \item Device Under Test (Quantity: 1)

    \item Resistor ($1 k\Omega$) (Quantity: 1)

    \item Multimeter (Quantity: 2)

\end{itemize}

\section{Procedure}
\begin{enumerate}

    \item Connect the circuit as per the experimental circuit diagram.

    \item Apply a constant gate current $I_G$.

    \item Vary the supply voltage from -15V to 15V.

    \item Record the voltage and current.

    \item Plot the characteristics.

\end{enumerate}

\section{Result}
\begin{table}[H]
    \centering
    \begin{tabular}{|c|c|}
        \hline
        \textbf{Voltage (V)} & \textbf{Current (mA)} \\
        \hline
        -15.0 & -759.3 \\
        -11.7 & -753.0 \\
        -8.4 & -7.6 \\
        -5.0 & -3484.0 \\
        -1.7 & -1.5 \\
        1.6 & 712.9 \\
        5.0 & 4.3 \\
        8.3 & 748.2 \\
        11.7 & 10.9 \\
        15.0 & 762.8 \\
        \hline
    \end{tabular}
    \caption{Simulated Data (Sample Points)}
\end{table}



\begin{figure}[H]
    \centering
    \includegraphics[width=0.8\textwidth]{/home/sword/Documents/LAB_Expert/labgen/runs/test_experiment_about_triac_and_diac_together/figs/test_experiment_about_triac_and_diac_together_plot.png}
    \caption{ Simulated Test Experiment about triac and diac together Curve }
\end{figure}



\section{Discussion}
In the laboratory setup, the DIAC was connected in series with the gate terminal of the TRIAC, while the TRIAC was connected in series with a resistive load to regulate AC power. The circuit was powered by a variable AC source, and the voltage across the load was monitored using an oscilloscope to observe the waveform changes. It was observed that the DIAC did not conduct until the source voltage reached its positive or negative breakover threshold, at which point it triggered the TRIAC into conduction. The resulting voltage waveform across the load exhibited phase control, with the conduction period starting later in each half-cycle as the firing angle increased. The measurement data indicated that the average power delivered to the load was reduced compared to the uncontrolled case, confirming the operation of the DIAC as a reliable trigger source for the TRIAC. The bidirectional nature of the components was verified, as the circuit functioned identically during both the positive and negative half-cycles of the input AC signal.

\section{Conclusion}
The experiment successfully demonstrated the operational principles of a DIAC-TRiac combination circuit for AC power control. The DIAC was observed to act as a bidirectional trigger, conducting only when the applied voltage exceeded its breakover point, which in turn forced the TRIAC into conduction. The results confirmed that the TRIAC maintains a latching state until the current commutates near the zero-crossing of the AC waveform, effectively enabling phase-controlled power delivery to the load. By varying the conditions that influence the DIAC's triggering time, the average voltage across the load was modified, illustrating the utility of these components in applications such as dimming and speed control. The observed waveforms aligned with theoretical predictions for phase-control circuits, validating the effectiveness of the DIAC in ensuring consistent triggering of the TRIAC across both polarity cycles. Future work could involve quantitative analysis of power factor and total harmonic distortion in the resulting load current.

\section{References}
\begin{enumerate}

    \item Generated by LabGen Knowledge Hub API

    \item Ngspice Simulation Data.

\end{enumerate}

\end{document}